% =====================================================================
%  4쪽 재구성본 컴파일 래퍼 — acmart sigconf
%  이전본 대비: 초록 제거, IRL 정식화 제거, 지연 자기행동/지속성 기준선 제거, 부록 제거.
%  템플릿 요구사항(00_템플릿_요구사항_정리.md): sigconf, CCSXML+ccsdesc, keywords,
%  모든 figure 에 \Description, booktabs, 세로줄 금지, ACM-Reference-Format.
%
%  컴파일 (paper-icaif26/ 에서):
%    cp build/build-ko-4p.tex . && xelatex build-ko-4p && bibtex build-ko-4p \
%      && xelatex build-ko-4p && xelatex build-ko-4p
% =====================================================================

% 제출본에서는 nonacm 을 제거하고, CFP 가 익명심사를 요구하면 anonymous 를 추가한다.
\documentclass[sigconf,nonacm]{acmart}

\IfFileExists{kotex.sty}{\usepackage{kotex}}{%
  \usepackage{fontspec}%
  \IfFontExistsTF{Noto Serif CJK KR}{%
    \setmainfont{Noto Serif CJK KR}[BoldFont={Noto Serif CJK KR Bold},
                                    ItalicFont={Noto Serif CJK KR},
                                    BoldItalicFont={Noto Serif CJK KR Bold}]%
    \setsansfont{Noto Sans CJK KR}[BoldFont={Noto Sans CJK KR Bold},
                                   ItalicFont={Noto Sans CJK KR},
                                   BoldItalicFont={Noto Sans CJK KR Bold}]%
    \setmonofont{Noto Sans Mono CJK KR}[Scale=0.92]%
  }{\setmainfont{Noto Serif CJK JP}\setsansfont{Noto Sans CJK JP}}%
  \XeTeXlinebreaklocale "ko"%
  \XeTeXlinebreakskip = 0pt plus 0.1em%
}

\usepackage{booktabs}
\usepackage{graphicx}
\usepackage{amsmath}
\usepackage{array}

% tabular 내부에서 행 조각 파일을 읽기 위한 원시 \input
\makeatletter
\newcommand{\inputrows}[1]{\@@input #1 }
\makeatother

\settopmatter{printacmref=false,printfolios=true}
\setcopyright{none}

\AtBeginDocument{%
  \renewcommand{\figurename}{그림}%
  \renewcommand{\tablename}{표}%
}

\begin{document}

\title{집계 수급에서 투자자 유형별 계수를 해석할 때의 식별 경계}
\subtitle{대칭 선호복원과 기계 성분 보정을 포함한 5단계 검증}

\author{저자명}
\affiliation{\institution{소속기관}\city{서울}\country{대한민국}}
\email{author@example.com}
\renewcommand{\shortauthors}{저자명}

\begin{CCSXML}
<ccs2012>
<concept>
<concept_id>10010405.10010455.10010461</concept_id>
<concept_desc>Applied computing~Economics</concept_desc>
<concept_significance>500</concept_significance>
</concept>
<concept>
<concept_id>10003752.10010070.10010071.10010261</concept_id>
<concept_desc>Theory of computation~Machine learning theory</concept_desc>
<concept_significance>300</concept_significance>
</concept>
</ccs2012>
\end{CCSXML}

\ccsdesc[500]{Applied computing~Economics}
\ccsdesc[300]{Theory of computation~Machine learning theory}

\keywords{투자자 유형별 수급, 벌점 회귀, 시장 청산 제약, 식별, 조합형 정화 교차검증}

\maketitle

% =====================================================================
%  [4쪽 재구성본 2026-07-27] SPR — 방법 / 실험 / 논의
%
%  이전본 대비 삭제: 역강화학습 정식화(보상함수·최적정책·clip·최대엔트로피 IRL),
%                    지연 자기행동 승격 사양과 지속성 기준선 일체, 초록.
%  추정량은 처음부터 벌점 회귀(Lasso)로 정의한다.
%  핵심 기여: (i) 5단계 검증 프로토콜, (ii) V5(기계 성분 보정)가 낳는 새 사실,
%             (iii) 집계 기관 범주의 행동적 비식별성.
%  의존: sections/generated/tab4p_body.tex, figures/fig4p_*.pdf
% =====================================================================

\section{방법: 대칭 선호복원}
\label{sec:method}

\textbf{기여.} 본 논문의 기여는 넷이다. \emph{첫째, 대칭 추정 설계.} 세 투자자 유형에
동일한 특징과 동일한 모형 형태를 부여하고 파라미터만 분리 학습하므로, 계수의 유형 간
차이는 설계가 아니라 자료에서 온다. \emph{둘째, 5단계 검증 프로토콜.} 복원된 계수를
그대로 보고하지 않고 표본 안정성 네 단계를 거치되, 여기에 다섯째 단계
\textbf{두 합성 대조}를 더한다. V5a는 아무도 특징에 반응하지 않는 세계로 파이프라인의
무선호 대역과 청산이 만드는 시차 교차연관을 구한다. V5b는 외국인에게 관측 반응을 부여하고
개인의 직접 반응만 0으로 둔 조건부 청산 세계를 만들어 개인 계수의 유도 대역을 구한다.
\emph{셋째, 두 대조가 낳는 실증 결과.} 외국인의 모멘텀ㆍ시장ㆍ환율 반응은 V5a 대역을
벗어난다. 개인의 모멘텀과 환율은 V5b 청산 대역 안이지만, 시장효과는 관측치 $-0.0238$이
조건부 대역 $[-0.0709,-0.0263]$보다 위에 있어 외국인 유도 매도를 일부 상쇄하는 양의
잔여 반응이 남는다. 음의 시차 교차연관은 V5a에서 같은 크기로 발생한다.
\emph{넷째, 집계 기관 범주의 행동적 비식별성.} 기관은 거래 규모가 가장 큼에도 주요
특징의 계수가 안정적이지 않고 표본외 설명력도 거의 0이다. 이는 기관이 정보를 고려하지
않는다는 뜻이 아니라, 여러 세부 주체를 합산한 기관 범주가 본 특징 공간에서 하나의
안정적인 공통 반응함수로 표현되지 않음을 보인다.

본 절은 대칭 선호복원(Symmetric Preference Recovery, 이하 SPR) 절차를 정의한다.

\subsection{행동 정의와 대칭 추정}
\label{sec:setup}

투자자 유형을 $\mathcal{I}=\{\text{외국인},\text{기관},\text{개인}\}$으로 둔다. 유형
$i$의 $t$일 총매수ㆍ총매도금액을 $V^{buy}_{i,t}$, $V^{sell}_{i,t}$라 하면 행동을
\begin{equation}
u_{i,t}=\frac{V^{buy}_{i,t}-V^{sell}_{i,t}}{V^{buy}_{i,t}+V^{sell}_{i,t}}\in[-1,1]
\label{eq:action}
\end{equation}
로 정의한다. 분모가 시장 전체가 아니라 해당 유형의 총거래금액이므로 $u_{i,t}$는 유형별
거래 규모에 영향받지 않는 공통 척도가 되며, 이진 라벨과 달리 강도 정보를 보존한다.

세 유형은 동일한 상태 특징 $x_{i,t}\in\mathbb{R}^{3}$, 동일한 컨텍스트
$C_t\in\mathbb{R}^{2}$, 동일한 모형 형태를 부여받고 파라미터
$\theta_i=(\beta_i,B_i,\alpha_i)$만 각자의 데이터로 분리 학습한다. 유형별로 다른 특징을
미리 배정하지 않으므로 결과의 유형 간 차이는 설계가 아니라 데이터에서 학습된 것이며,
이것이 계수의 직접 비교를 정당화한다. 학습 표본은 $(x_{i,t},C_t,a_{i,t+1})$,
$a_{i,t+1}=u_{i,t+1}$이다. $t$일 상태가 $t+1$일 행동을 설명하므로 미래 거래는 특징에
들어가지 않는다.

\subsection{상태 특징과 시장 컨텍스트}
\label{sec:features}

세 특징과 두 컨텍스트는 모두 각 분할의 학습 인덱스에서만 표준화하므로, 이후 계수는 1
표준편차 변화의 효과로 읽는다.

\textbf{(F1) 중기 모멘텀} $x^{mom}_{t}=\log(P_t/P_{t-20})$. 20거래일은 해석을 위한 사전
지정값이며 자료에서 최적화하지 않았다. 양수는 추세추종, 음수는 역추세를 뜻하고 부호
기대는 유형별 모멘텀 거래
문헌~\cite{grinblatt2000investment,choe1999foreign}에서 온다.

\textbf{(F2) 유형 간 시차 교차연관} 유형 $i$를 제외한 두 유형의 전일 행동 횡단면 평균
\begin{equation}
x^{cross}_{i,t}=\tfrac{1}{2}\left(u_{j,t-1}+u_{k,t-1}\right),\qquad j,k\neq i.
\label{eq:cross}
\end{equation}
1거래일 지연은 정보 제약이 아니라, 동일 거래일 유형 간 수급의 회계적 상쇄를 예측
변수에서 배제하기 위한 설계 선택이다. 부호 기대는 기관 군집행동
문헌~\cite{lakonishok1992impact}을 참조하되, 해당 문헌은 기관 \emph{내부}의 군집을,
식~\eqref{eq:cross}은 \emph{유형 간} 시차 관계를 측정하므로 이 계수는 군집행동의 직접
검정이 아니다.

\textbf{(F3) 평균단가 대비 손실구간} 평균단가 대용치 $\bar P_{i,t}$와 당일 총거래 VWAP
$P^{vwap}_{i,t}$로 $x^{uw}_{i,t}=\max\{(\bar P_{i,t}-P^{vwap}_{i,t})/\bar P_{i,t},\,0\}$.
0에서 절단하므로 준거점 의존적 설계다. 평균단가는 총매수ㆍ총매도 수량과 금액으로 유효
보유량을 재귀 갱신해 구하되 전일 보유량에 망각계수 $\rho=0.98$을 곱한다(반감기 약
34거래일). 부호 기대는 처분효과
문헌~\cite{shefrin1985disposition,odean1998investors}에서 오나, 이 대용치는 집계 거래로
만든 근사치이므로 계수를 처분효과의 직접 검정으로 읽을 수 없다.

\textbf{컨텍스트}는 KOSPI200 1일 수익률 $r_t$와 원/달러 환율 수준
$z_t=(\log e_t-\mu_{t-1,252})/\sigma_{t-1,252}$이다. 평균과 표준편차는 $t-1$까지의
252거래일만 쓴다. 삼성전자의 KOSPI200 비중 때문에 종목 모멘텀과 시장수익률이 겹치며
특징--컨텍스트 상관은 $-0.29$에서 $0.22$ 범위다. 따라서 개별 계수는 다른 항이 통제된
조건부 연관이다.

\subsection{추정: 벌점 회귀와 해의 유일성}
\label{sec:model}

유형 $i$의 행동점수를 $q_{i,t}=(\beta_i+B_iC_t)^{\top}x_{i,t}+\alpha_i^{\top}C_t$로 둔다.
$\alpha$는 종목 특징이 평균일 때도 행동 수준을 옮기는 수준효과, $B$는 컨텍스트에 따라
특징 민감도가 달라지는 기울기 효과다. 상호작용을 넣으면서 하위차수 항 $C_t$를 빼는 것은
표준적 사양 오류이므로 $\alpha$의 포함은 자료로 정당화할 대상이 아니다.
$z_{i,t}=[\,x_{i,t};\,\mathrm{vec}(x_{i,t}C_t^{\top});\,C_t\,]\in\mathbb{R}^{11}$,
$\theta_i=[\,\beta_i;\,\mathrm{vec}(B_i);\,\alpha_i\,]$로 두면 $q_{i,t}=\theta_i^{\top}z_{i,t}$
이고, 추정량은 학습 인덱스 $\mathcal{D}_i$ 위의 벌점 회귀다.
\begin{equation}
\hat\theta_i=\operatorname*{arg\,min}_{\theta}\ \Bigl\{\tfrac{1}{|\mathcal{D}_i|}
\!\!\sum_{t\in\mathcal{D}_i}\!\!\bigl(z_{i,t}^{\top}\theta-a_{i,t+1}\bigr)^{2}
+\lambda_i\lVert\theta\rVert_1\Bigr\}
\label{eq:loss}
\end{equation}
유형당 파라미터는 11개($\beta$ 3, $B$ 6, $\alpha$ 2)다. 식~\eqref{eq:loss}은 볼록이므로
좌표하강으로 전역해를 정확히 구한다. 다만 볼록성만으로 해가 유일한 것은 아니어서 두
조건을 분할마다 확인한다. $\lambda_i=0$이면 $\mathrm{rank}(Z_i)=11$일 때 목적식이
강볼록이라 유일하고, $\lambda_i>0$이면 등상관 집합
$E=\{j:|\tfrac{2}{n}z_j^{\top}(a-Z\hat\theta)|=\lambda_i\}$에 대해
$\mathrm{rank}(Z_{i,E})=|E|$일 때 유일하다. 45개 분할$\times$3유형 전부에서 두 조건이
성립했다(최악 조건수 6.86, $\lambda>0$인 두 유형의 $|E|=8$ㆍ$11$). 즉 보고하는 계수는
각 분할에서 유일하게 결정된 값이다.

\subsection{검증 프로토콜}
\label{sec:protocol}

복원된 계수를 그대로 보고하지 않고 다섯 단계를 거친다. 앞 넷은 표본 안정성을, 마지막
하나는 식별을 다룬다.

\textbf{(V1) 조합형 정화 교차검증.} 10개 시간 폴드 중 2개를 시험으로 쓰는
CPCV~\cite{lopez2018advances}로 $\binom{10}{2}=45$개 분할을 만들고, 1일 purge와 5일
embargo를 적용하며 모든 표준화는 학습 인덱스에만 적합한다. 분할들이 관측을 공유하므로
부호 일관성은 $p$-값이 아니라 안정성 지표로만 해석한다.

\textbf{(V2) 달력월 블록 부트스트랩.} 그 한계를 보완하기 위해 표본을 달력월 블록으로
나누어 1{,}000회 복원추출하고 95\% 신뢰구간을 구한다. 월 블록은 일별 수급의 자기상관을
블록 안에 보존한다.

\textbf{(V3) Expanding walk-forward.} CPCV는 시험 폴드가 학습 폴드보다 앞설 수 있으므로,
최초 400개 관측에서 시작해 5일 purge 후 다음 60거래일을 예측하고 학습 구간을 60거래일씩
넓히는 창 10개를 수행한다.

\textbf{(V4) 벌점 교체.} 결과가 $L_1$ 벌점의 산물이 아님을 보이기 위해 $L_2$ 및
무벌점으로 재추정해 부호를 비교한다.

\textbf{(V5) 두 합성 대조.} 앞의 네 단계는 계수가 \emph{안정적}인지만 묻는다. V5a는
모든 유형의 직접 반응을 0으로 둔 무선호 합성 수급으로 파이프라인 영가설을 보정한다.
V5b는 외국인의 관측 반응을 유지하고 개인의 직접 반응만 0으로 둔 뒤 관측 흡수율로
청산시켜, 개인 계수의 조건부 청산 대역을 구한다(\S\ref{sec:v5}). 같은 추정 파이프라인을
쓰되 서로 다른 귀무를 검정한다.

지표는 방향 정확도($\mathrm{sign}(\hat a)=\mathrm{sign}(a)$의 비율), 예측--실제 상관,
MAE다. 유형 간 행동 표준편차가 다르므로 오차의 절대값은 유형 간 비교에 쓰지 않는다.

% =====================================================================
\section{실험}
\label{sec:experiments}

\subsection{데이터}
\label{sec:data}

삼성전자(005930) 일별 자료에서 252일 환율 이력, 20일 모멘텀, 다음 날 행동을 모두 계산할
수 있는 최신 973행을 사용한다. 상태일은 2022-01-06--2025-12-29이다. 단일 종목은 표본의
한계가 아니라 종목 고정효과와 유동성 차이를 배제하고 유형 간 이질성만 분리하기 위한
통제된 설정이다.

행동 평균은 세 유형 모두 0에 가깝지만 표준편차는 개인 0.334, 외국인 0.257, 기관 0.125로
다르고 1차 자기상관은 각각 0.345, 0.401, 0.149다. 동일 거래일 유형 간 행동은 강한 음의
관계를 갖는다(외국인--개인 $-0.772$, 기관--개인 $-0.499$). 세 유형 순매수 합계와 시장
전체의 차이(기타법인ㆍ기타외국인 등)는 일평균 144억 원, 거래대금의 1.3\%에 그친다.
$\lambda$는 유형별로 $0$ㆍ$0.01$ㆍ$0.0003$이며, 이는 본 사양에서 중첩 교차검증으로 고른
값이 아니라 선행 사양에서 이월된 값이다. 세 유형에 동일한
$\lambda\in\{0,0.0005,0.01\}$을 부여해도 아래의 부호 판정은 바뀌지 않는다.

\subsection{계수와 표본 안정성 (V1--V4)}
\label{sec:stability}

\begin{figure*}[t]
  \centering
  \includegraphics[width=\textwidth]{figures/fig4p_forest.pdf}
  \caption{복원된 계수의 안정성과 두 합성 대조. 회색 점은 45개 개별 CPCV 분할, 굵은 선은
  분할 간 표준편차, 옅은 선은 달력월 블록 부트스트랩 95\% 신뢰구간이다. 속이 빈 표식은
  신뢰구간이 0을 포함하는 항이다. 회색 띠는 V5a 무선호 95\% 대역, 보라 띠는 개인
  모멘텀ㆍ두 컨텍스트에 대한 V5b 조건부 청산 95\% 대역이며, 오른쪽 숫자는 45개 분할의
  평균 계수다.}
  \Description{두 패널의 forest plot. 왼쪽은 모멘텀, 시차 교차연관, 손실구간 계수를,
  오른쪽은 KOSPI200과 환율 주효과를 세 투자자 유형에 대해 보인다. 개인 모멘텀과 환율은
  보라 조건부 대역 안, 개인 KOSPI200 효과는 그 대역의 위쪽 밖에 놓인다.}
  \label{fig:forest}
\end{figure*}

\begin{table}[t]
  \caption{CPCV 표본외 성능(45개 분할 평균 $\pm$ 표준편차). 예측 성능의 우위를 주장하기
  위한 표가 아니라 모형이 잡음만 학습한 것이 아님을 보이는 표다.}
  \label{tab:performance}
  \small
  \centering
  \begin{tabular}{@{}lccc@{}}
    \toprule
    지표 & 외국인 & 기관 & 개인 \\
    \midrule
    방향 정확도 & $0.636\pm0.040$ & $0.536\pm0.035$ & $0.602\pm0.044$ \\
    상관계수    & $0.279\pm0.106$ & $0.060\pm0.062$ & $0.239\pm0.116$ \\
    MAE         & $0.195$ & $0.095$ & $0.266$ \\
    표본외 $R^2$ & $0.105$ & $0.002$ & $0.102$ \\
    \bottomrule
  \end{tabular}
\end{table}

\begin{figure*}[t]
  \centering
  \includegraphics[width=\textwidth]{figures/fig4p_prediction.pdf}
  \caption{투자자별 실제 행동과 CPCV 표본외 예측 행동. 각 날짜의 9개 CPCV 시험예측을
  먼저 평균한 뒤 실제값과 예측값에 20거래일 이동평균을 적용했다. 패널 제목의
  표본외 $R^2$는 평활 전 일별 예측에 대한 45개 분할 평균으로
  표~\ref{tab:performance}와 같다.}
  \Description{외국인, 기관, 개인의 실제 순매수 행동과 표본외 예측을 시간축 위에
  겹쳐 그린 세 개의 선그래프. 외국인과 개인의 예측은 실제 행동의 완만한 국면을
  따라가지만, 기관의 예측은 0선 부근에 거의 평탄하게 놓여 있다.}
  \label{fig:prediction}
\end{figure*}

표~\ref{tab:performance}에서 분할 간 산포가 커서 단일 분할의 수치로 유형을 비교할 수
없다. 그림~\ref{fig:prediction}에서는 외국인과 개인의 예측이 완만한 수급 국면을
따라가지만 기관의 예측은 0선 부근에 붙어 있으며, 계수 안정성에서도 같은 대비가
나타난다(그림~\ref{fig:forest}).

\textbf{(V1)} 외국인의 모멘텀 가중치는 $+0.0502$로 45개 분할 전부에서 양수, 개인은
$-0.0383$으로 전부에서 음수다. 시차 교차연관은 세 유형 모두 전 분할에서 음수이고, 개인의
손실구간은 97.8\%의 분할에서 양수다. 기관은 모멘텀과 손실구간이 45개 분할 중 각각 36개,
44개에서 정확히 0으로 소거된다.

\textbf{(V2)} 부트스트랩에서 0을 제외하는 항은 8개다. 외국인의 모멘텀ㆍ시차
교차연관ㆍ$\alpha^{KOSPI}$ㆍ$\alpha^{FX}$, 개인의 같은 네 항이며, 개인 손실구간
($[-0.0122,+0.0665]$)과 기관의 모든 항은 통과하지 못한다. 개인 손실구간은 45개 분할에서
부호가 100\% 일관됨에도 월 블록 재표집의 5.5\%가 부호를 뒤집는데, 이는 CPCV 분할이 관측을
공유하기 때문에 부호 일관성이 불확실성을 과소평가한다는 \S\ref{sec:protocol}의 경고가
실제로 작동한 사례다.

\textbf{(V3)} Expanding walk-forward에서 성능은 창마다 흔들린다(외국인 방향 정확도
$0.610\pm0.116$, 상관 $0.131\pm0.162$로 CPCV의 절반 이하). 본 연구는 이 저하를 시간
구조에서 오는 일반화 한계로 그대로 보고한다. 반면 계수 경로는 흔들리지 않는다. 외국인
모멘텀은 10개 창 전부 양수, 개인은 전부 음수이며 두 경로는 한 번도 교차하지 않고, 두
컨텍스트 주효과에서도 같은 대비가 유지된다. 기관의 경로는 0선에 붙어 있다.

\textbf{(V4)} 외국인과 개인의 계수는 $L_1$, $L_2$, 무벌점에서 소수 넷째 자리까지 일치한다.
기관의 계수만 벌점에 따라 이동한다(시차 교차연관 $-0.0076\rightarrow-0.0124$).

\subsection{두 합성 대조 (V5)와 그 귀결}
\label{sec:v5}

\begin{table*}[t]
  \caption{안정 계수의 두 합성 대조. 계수는 정확해 CPCV 45개 분할 평균, 괄호는 부호
  일관성(\%)이다. V5a는 모두 무선호인 귀무, V5b는 외국인 관측 반응과 개인 직접 반응 0인
  조건부 청산 귀무의 1{,}000회 95\% 대역이다.}
  \label{tab:filters}
  \small
  \setlength{\tabcolsep}{4pt}
  \centering
  \begin{tabular}{@{}llccccc l@{}}
    \toprule
    항 & 유형 & 계수 (부호 일관성) & 부트스트랩 CI가 0 제외 & 대조
      & 95\% 대역 & 대역 밖 & 판정 \\
    \midrule
    \inputrows{sections/generated/tab4p_body.tex}
    \bottomrule
  \end{tabular}
\end{table*}

\textbf{청산 제약이 유도하는 성분.} 유형 간 계수를 비교하기 전에, 시장 청산이 기계적으로
만들 수 있는 것부터 계산한다. 본 표본에서 외국인 순매수 1원은 평균적으로 개인 $-0.92$원,
기관 $-0.06$원, 기타 $-0.02$원으로 흡수되고 유형별 총거래금액 비율은
$\mathbb{E}[W_f/W_r]=1.21$--$1.42$다. 따라서 외국인이 표준화 특징 $x$에 $\beta_f$만큼
반응하면 개인이 아무 선호가 없어도 개인 계수에는
$\beta^{ind}_{r}\approx-0.92\times\beta_f\times(1.21\text{--}1.42)$의 성분이 유도된다.
관측된 외국인 계수를 대입하면 유도 성분은 모멘텀 $-0.056$--$-0.066$,
$\alpha^{KOSPI}$ $-0.040$--$-0.047$, $\alpha^{FX}$ $+0.030$--$+0.035$이고, 이는 관측된 개인
계수($-0.0383$, $-0.0238$, $+0.0355$)의 각각 146--172\%, 168--198\%, 84--99\%에 해당한다.
관측치는 이 예측의 각각 58--68\%, 51--60\%, 101--118\%이므로, 단순히 ``가능 범위 안''인지
아니라 순수 흡수자 귀무에서 체계적으로 벗어나는지를 조건부로 검정해야 한다.

\textbf{V5a: 무선호 합성 대조.} 외국인ㆍ기관 행동과 기타 잔차를 관측 표준편차ㆍ자기상관에
맞춘 특징 비관련 잡음으로 만들고 개인을 청산시킨다. 모멘텀과 컨텍스트 대역은 0 중심의
파이프라인 귀무분포다. 반면 시차 교차연관은 \emph{무선호에서도} 음수가 되며(외국인
$[-0.057,-0.018]$, 개인 $[-0.085,-0.011]$), 관측치 $-0.0406$ㆍ$-0.0466$은 내부에 있다.
따라서 이 항의 음수는 청산과 자기상관만으로 발생한다.

\textbf{V5b: 조건부 청산 대조.} 외국인에게 관측된 모멘텀ㆍ시장ㆍ환율 반응을 부여하고
개인의 직접 반응은 0으로 둔다. 외국인 순매수의 흡수율(기관 $-0.059$, 개인 $-0.924$,
기타 $-0.017$)로 신호를 배분한다. 네 주체의 특징 비관련 금액잔차는 공동 원형 이동하여
특징과의 정렬만 끊고, 모든 잔차에 공통 척도 0.962를 적용한다. 이 방식은 일별 청산식과
잔차의 시간순서ㆍ횡단면 관계를 보존한다. 합성 개인의 표준편차 $0.3338$, 외국인 상관
$-0.770$, 기타 수급 표준편차 222억 원은 관측 $0.3339$, $-0.772$, 230억 원과 가깝다.

\textbf{조건부 판정.} 개인 모멘텀의 V5b 대역은 $[-0.0963,-0.0242]$으로 관측 $-0.0383$을,
환율 대역은 $[-0.0080,+0.0786]$로 관측 $+0.0355$를 포함한다. 두 부호는 개인 고유 선호
없이 설명된다. 그러나 개인 시장효과의 대역은 $[-0.0709,-0.0263]$이고 관측 $-0.0238$은
상단 밖(98.4 백분위)이다. 관측계수는 여전히 음수지만 조건부 평균 $-0.0494$보다
$+0.0256$만큼 크다. 즉 개인에게는 외국인 유도 매도를 일부 상쇄하는 양의 잔여 시장반응이
남는다. 외국인의 모멘텀ㆍ시장ㆍ환율은 V5a 무선호 대역을 벗어나 고유 반응으로 남는다.

\textbf{강건성.} 망각계수 $\rho\in\{0.95,0.98,0.99\}$와 모멘텀 창
$\{10,20,60\}$의 아홉 조합에서 외국인ㆍ개인 모멘텀은 각각 100\% 양ㆍ음이고,
개인 손실구간은 91\% 이상 양수다. 세 유형에 $\lambda$를 통일해도 같다. 합성 대조의 보정 한계는 그대로
보고한다. V5b는 관측 흡수율이 특징이 유발한 거래에도 같다는 동질 흡수 가정에 의존하므로,
조건부 대역은 구조적 진실이 아니라 명시된 귀무 아래의 식별 기준이다.

% =====================================================================
\section{결과 해석과 결론}
\label{sec:interpretation}

\textbf{외국인: 시장 방향과 환율 국면에 대한 독자적 반응.} 모멘텀 가중치
$+0.0502$는 45개 CPCV 분할에서 모두 양수이고 부트스트랩 신뢰구간이 0을 제외하며,
기계 대역도 $4.3\sigma$ 벗어난다. 따라서 최근 주가 상승 뒤 외국인 순매수가 증가하는
관계는 거래상대방 구조만으로 유도된 것이 아니라 외국인 수급에 남은 안정적 추세추종
반응으로 읽힌다. KOSPI200 효과 $+0.0360$과 환율 효과 $-0.0269$도 두 필터를 통과한다.
즉 시장 상승일에는 순매수 쪽으로, 원/달러 환율이 직전 252일 기준으로 높은 원화 약세
국면에는 순매도 쪽으로 이동하며, 외국인 추세거래 문헌의 부호 기대와
일치한다~\cite{grinblatt2000investment,choe1999foreign}. 반면 음의 시차 교차연관은
기계 대역 안에 있으므로 역군집 성향으로 해석하지 않는다.

\textbf{개인: 청산의 반대편과 제한적 잔여 반응.} 개인의 모멘텀 $-0.0383$과 환율
$+0.0355$는 안정적이지만 V5b 대역 안이므로 외국인 거래의 흡수만으로 설명된다. 시장효과
$-0.0238$은 조건부 청산 대역보다 덜 음수다. 따라서 개인이 시장 상승일에 순매수한다는
뜻은 아니지만, 외국인 유도 매도압력을 일부 상쇄하는 양의 잔여 반응은 식별된다. 손실구간
가중치는 $+0.0315$이고 CPCV 분할의 97.8\%와 강건성 사양의 91\% 이상에서 양수다.
이는 평균단가 대비 손실이 깊어질수록 순매수하거나 매도를 줄이는 방향으로, 손실실현 회피
또는 추가매수와 부호상 정합한다. 다만 월 블록 신뢰구간이 0을 포함하고, 집계 순매수와
평균단가 대용치로는 추가매수와 매도 보류를 구분할 수 없으므로 손실회피의 직접 증거가
아닌 반복적 행동 신호로 한정한다.

\textbf{기관: 집계 수급에서 나타나지 않는 공통 행동가중치.} 개별 기관이 어떤 정보를
사용하는지는 알 수 없지만, 한국 시장의 집계 ``기관'' 수급에서는 본 연구가 설정한
행동특징들이 하나의 안정적인 공통 매매 원리로 나타나지 않는다. 모든 주요 계수가
불안정하고 표본외 $R^2$가 0.002인 결과는 세부 기관의 이질적 반응이 합산 과정에서
상쇄될 가능성과 정합한다.

\textbf{종합.} 외국인의 세 반응은 무선호 대조 뒤에도 남는다. 개인 모멘텀ㆍ환율은 청산의
반대편이지만 시장효과에는 양의 잔여 반응이 있고, 손실구간에는 손실실현 회피와 정합적인
제한적 신호가 남는다. 기관에서는 공통 행동원리가 식별되지 않는다. 계수의 안정성만이
아니라 검정하려는 귀무에 맞춘 합성 대역과 함께 해석해야 한다는 것이 결론이다.

% =====================================================================
\section{한계와 향후 과제}
\label{sec:limitations}

\textbf{(1) 해석 범위.} 복원된 가중치는 인과효과가 아닌 조건부 연관이며 수익률ㆍ샤프지수ㆍ
거래비용을 포함하지 않는다. 따라서 표~\ref{tab:performance}에서 매매 규칙이나 수급 예측의
우위를 도출하지 않는다. 시차 교차연관은 청산이 유도하는 성분과 분리되지 않으며, 개인
V5b의 잔여 반응은 동질 흡수율과 합성 잔차모형을 조건으로 한다.

\textbf{(2) 처분효과 대용치.} 개인 손실구간은 월 블록 부트스트랩을 통과하지 못했고,
평균단가는 계좌별 실현손익이 아니라 집계 거래로 만든 근사치다. 따라서 처분효과의 직접
검정이 아니라 부호 정합의 관측이다.

\textbf{(3) 일반화.} 단일 종목 설정은 유형 간 이질성 분리를 위한 통제이나 다른 종목과
시장으로의 확장은 검증되지 않았다.

\textbf{향후 과제.} 첫째, 기계 대역의 폭은 종목별
투자자 구성비 $\mathbb{E}[W_i/W_j]$와 흡수 몫의 함수이므로, 외국인 지분과 유동성이 다른
종목에서는 대역 대비 관측 계수의 위치가 예측 가능하게 달라져야 한다. 이것이 기계적 성분과
선호 성분을 종목 횡단면에서 분리하는 설계다. 둘째, 기관의 null이 합산에서 온다는 가설은
KRX 세부 투자자별 자료로 직접 검정할 수 있으며, 본 논문의 대칭 추정과 V1--V5는 세부 주체
수에 관계없이 그대로 적용된다.


\bibliographystyle{ACM-Reference-Format}
\bibliography{paper-ko,paper-refs-additional}

\end{document}
